\textbf{Date:} November 17, 2026
\\\textbf{Interviewer:} Lucas DeBoer
\\\textbf{Participant:} P3, a nearly life-long quadriplegic who still has their vision, but controls many things in his day to day life using voice controls.
\\\textbf{Format:} In-person conversation, about 30 minutes long
\\\textbf{Editing:} The entire conversation was transcribed and recorded using Apple's Voice Memo app. Light editing was done afterwards to fix transcription mistakes and keep the transcript readable by keeping Lucas' one-word replies ("Yeah," "Yep," "Gotcha") inside P3's turns as \textit{[Lucas: Yeah.]}. Personally identifying information and information not relevant to the interview were also removed.
\\\textbf{Recruitment:} A neighbour of the interviewer, recruited directly
\\\textbf{Consent:} P3 verbally consented to the interview and its recording prior to the recording beginning.
\\\textbf{Anonymization:} P3 is not named. All identifying details have been removed.

\bigskip
\hrule
\bigskip

\textbf{Lucas:} Okay, so, I gave the intro. So we're just learning about what works well for you or things that are frustrating, kind of how you use the computer. Obviously, this isn't a test, so there's no right or wrong answers. We're just interested in your experiences, so you can skip any questions, stop, whatever you like. 

\textbf{P3:} Okay, yeah. 

\textbf{Lucas:} Um, so if you could tell me a little bit about yourself and about how technology fits into your day, just at a high level, as much as you're comfortable sharing, how does your disability affect how you use a computer? 

\textbf{P3:} Um, well, I, I did have my hands. I was in a motorcycle accident. I'm a quadriplegic, and I had been for 50 years, believe it or not, time flies. And I use, I need it for everything. I had my arms up until two years ago, and then I lost the use of my arms, just from the original accident. \textit{[Lucas: Yeah.]} And this time, I guess. And I really rely on it now for my voice. I used to drive with my arms, my wheelchair. \textit{[Lucas: Yeah.]} Now, it's with my head. I have to drive with my head. My computer, my television, everything I used to use, remote control, for, or, uh, splints that would go on my hands, so I could use them. And now I have to use voice for all of that. 

\textbf{Lucas:} Gotcha. Yeah. 

\textbf{P3:} Now, they have made, uh, a lot of improvements just a couple years since I started to use my voice, but I'm still limited to what I can do. I can't do as much as I used to, but it's pretty amazing what they have, yeah. 

\textbf{Lucas:} Yeah. Yeah. So what are some tasks that you do use your computer for? 

\textbf{P3:} I used to use it for writing, paying bills, like, for banking. \textit{[Lucas: Yeah.]} And I really like that. I'm, that's one thing I'm quite frustrated with, but I understand. My banking, for paying bills, for doing any investments for that. My easiest route, what I use the computer for, I can use my voice commands. \textit{[Lucas: Yeah.]} However, when I'm in banking, it won't let me use voice commands. 

\textbf{Lucas:} Okay. 

\textbf{P3:} So I have a thing that's there for my computer that I use my mouth with. \textit{[Lucas: Okay.]} to move with, it's my mouse. 

\textbf{Lucas:} Gotcha. 

\textbf{P3:} It's called a LipSync. 

\textbf{Lucas:} Okay. 

\textbf{P3:} And I believe it's Canadian. 

\textbf{Lucas:} Oh. 

\textbf{P3:} Because they got it for everything from British Columbia. 

\textbf{Lucas:} Okay, yeah. 

\textbf{P3:} So that, that's truly excellent. And they, they're incredible. Like, I've used it. They've had to replace it likely four or five times with updates, or because I'm having difficulty using it anymore. \textit{[Lucas: Yeah.]} And, uh, zero. They charge me nothing. 

\textbf{Lucas:} That's awesome. 

\textbf{P3:} I get it, yeah, it's incredible. It's a really great service, and, uh, yeah, good to be part of it, but that's one issue, and the reason is the bank is, uh... rightfully concerned about, uh... Uh, what am I thinking of? They don't want anybody getting into it. \textit{[Lucas: Yeah.]} Secure your account or into their banking system. So it's, I, security is the word I'm looking for. \textit{[Lucas: Yeah.]} So I can't use the voice commands when I'm in banking at all. 

\textbf{Lucas:} Gotcha. 

\textbf{P3:} It just won't let me. And it's frustrating. Because I'll get on, and that's the other thing about it, it sometimes takes longer. And so I get cut off. I spend 10 minutes going onto a banking site, and that tells me that I'm all timed out \textit{[Lucas: Gotcha.]}, and I have to do it all again. So that's very frustrating. 

\textbf{Lucas:} Gotcha. So those are two of the main things. 

\textbf{P3:} Yeah. Just the time, uh, I'm not quick enough to get things completed in time, and that, for security. \textit{[Lucas: Yeah.]} I can't use it there. So that's, uh, that's a real nuisance. 

\textbf{Lucas:} Gotcha. So you mentioned LipSync. Are there any other accessibility tools, features, or voice assistance that you use? 

\textbf{P3:} No, there's not. There are programs out there. \textit{[Lucas: Yeah.]} Uh, but I don't ..., the voice one I use is called Voice Access. 

\textbf{Lucas:} Voice Access, okay. 

\textbf{P3:} And it's, it's really good. It's getting better because there are still limitations to what you can, you can't do. \textit{[Lucas: Yeah.]} But, like, with the word processing, I can still write, I can cut and paste, but there's a lot of the words, you have to pronounce everything just perfectly. \textit{[Lucas: Yeah.]} Or you get something silly. \textit{[Lucas: Yeah.]} That's quite frustrating about it too. \textit{[Lucas: Yeah.]} It's good, but it's a long ways to go because if it's similar ... It's funny, you get that, so these things. I say one thing and I look up and its not even remotely close to what I'm talking about. \textit{[Lucas: Yeah. Yeah.]} 

\textbf{Lucas:} Um, so... The space is improving quickly. Like, the tools are getting better. 

\textbf{P3:} Yes, yeah, definitely are. 

\textbf{Lucas:} Okay. So you've kind of mentioned that you're using voice commands. What exactly do you mean by that? Like, if you're trying to, say you're trying to open up Word and type something, when you say you're using voice commands, can you explain that a little more? 

\textbf{P3:} Okay, yeah, I can... I have, the program is called Voice Access. I have to have, if I've got it running in the background, I have to say, "Voice Access, wake up." 

\textbf{Lucas:} Okay. You have to wake it up. 

\textbf{P3:} Yeah. And then you have to say something, yeah, they have to mute it at the end to shut it down. \textit{[Lucas: Yeah.]} Or it would copy down every word you and I are saying. \textit{[Lucas: Yeah.]} If I didn't have it. So there's different commands like that. But in word processing, yeah, I can copy, and cut and paste, and basic commands like that. \textit{[Lucas: Yeah.]} Fairly easy, but some of the more complex commands, it's not too good at, but they're improving. \textit{[Lucas: Yeah.]} They're improving. 

\textbf{Lucas:} So what would be an example of a complex command like that? 

\textbf{P3:} Hmm. Uh, in my spreadsheet, copying formulas and things like that. 

\textbf{Lucas:} Gotcha. 

\textbf{P3:} It makes me a little bit nervous because if you say one thing or you copy, it can delete your your whole page of work or your file you have. 

\textbf{Lucas:} Okay, so you've mentioned banking. Is there any other times that you've had difficulty completing something on a computer? What were you trying to do and what happened? 

\textbf{P3:} Yeah, let me think for a sec. 

\textbf{Lucas:} Yeah, no worries. 

\textbf{P3:} It's mainly, uh, with spreadsheets. 

\textbf{Lucas:} Okay. 

\textbf{P3:} I worry about that because you put in a formula, or you put in a number, and then you do one thing or say one thing, and it'll delete whatever. \textit{[Lucas: Yeah.]} And then you gotta start it all over again. 

\textbf{Lucas:} Gotcha. 

\textbf{P3:} So that's, yeah. They've got a lot of work to do to get that correct for everyone. 

\textbf{Lucas:} Yeah. So your voice command, it's not that it's running on its own. It's not like you say, open up spreadsheets and make me a graph. You're telling it explicitly, copy this cell, paste to this cell. 

\textbf{P3:} Right, okay. Yeah. 

\textbf{Lucas:} Right. So it's not really running on its own... 

\textbf{P3:} And I understand the difficulty from their side because they set it up so that it will work. But you really have to train it to your voice, and even though you do that, it will still, uh, do commands that you don't want. \textit{[Lucas: Yeah.]} It's, it's accuracy, I guess. Because you can spend an hour on a spreadsheet, then say one thing and think you said something else, \textit{[Lucas: Yeah.]} and bang. All your work's gone. 

\textbf{P3:} Yeah, exactly. 

\textbf{Lucas:} Um, so, if you need help with a computer task, say there's something that you can't do, like, say your banking doesn't work, what do you usually do? Like, what are your options? 

\textbf{P3:} Oh, I have to try to, I have really two things, two choices. I either have to get help from somebody who can physically use the keyboard to help me, which I don't like to do, 'cause, well, its in banking or whatever. \textit{[Lucas: Yeah.]} Unless it's my wife, I don't really want them to see it, or anything like that. \textit{[Lucas: Yeah.]} 

\textbf{Lucas:} Um, but you are comfortable letting, if it could do it, you would become comfortable letting your Voice Access into your bank? 

\textbf{P3:} I am. I am, but I am nervous about all this AI stuff. \textit{[Lucas: Yeah.]} Because I rely, not just on the computer, like, but you have to use my voice, which, security wise, is not as safe. 

\textbf{Lucas:} Mmmm, Yes. 

\textbf{P3:} So that's definitely a concern. \textit{[Lucas: Yeah.]} And I just, like, even just with my banking, uh, like, hackers can get into the American military. \textit{[Lucas: Yeah.]} So, I mean, they could get into mine and empty my bank account. \textit{[Lucas: Yeah.]} So that's definitely a concern with AI that worries me. 

\textbf{Conversation about AI...} 

\textbf{P3:} No, I, uh, I would like to avoid it altogether. But the other day, I either read or saw on a program on TV that night. I don't remember for sure which. But it was about... Where was they going with this? I lost my train of thought there. \textit{[Lucas: Yeah.]} Sorry, I have to come back. 

\textbf{Lucas:} That's fine. You were talking about AI, how you wanted to avoid it, and you had seen some article or something. 

\textbf{P3:} Yeah. Yeah. I wanted to avoid AI, so I am avoiding it, I thought, pretty well completely. \textit{[Lucas: Yeah.]} And then I heard or read that if you use Google, you're using AI. 

\textbf{Lucas:} Yeah. That blurb right at the top. 

\textbf{P3}: So there's no way really getting around it. \textit{[Lucas: Yeah.]} And it's got a lot of good features about it that, you know, help people like myself \textit{[Lucas: Yeah.]} and able-bodied people. It's like, wait, it's amazing how much it helps, but it's just, uh, I guess it's a whole security issue. \textit{[Lucas: Yeah.]} That you got to worry about. And also environmentally, they're talking about these data centers. That you just, a ridiculous amount of energy and water. \textit{[Lucas: Yeah.]} And yeah, so that's a, that's a concern. There's one, you know, if you jump on the 403 to go to Toronto, and, uh, before you get on the QEW. \textit{[Lucas: Yep.]} There's that one area there. They're going to put a data centre in there. \textit{[Lucas: Yeah.]} If they get permission. 

\textbf{ More conversation about AI...} 

\textbf{P3:} Um... Another thing I really use voice for a lot. \textit{[Lucas: Yeah.]} I don't know whether it will tie to this at all or not. I have my television. \textit{[Lucas: Okay.]} But for changing channels and, uh, volume and every other feature that I have on television, that's really good. They've got it fine-tuned compared to the computers. As far as I'm concerned. 

\textbf{Lucas:} Oh, okay. Yeah. So is that a, is that a separate? Is that, like, through your TV remote? Is that a separate program? 

\textbf{P3:} I think it's a... Some of it's the same. \textit{[Lucas: Okay.]} Some of it's Google, some of it's that Voice Access. Uh, yeah, it's, uh... \textit{[Lucas: Gotcha. Yep.]} They're the main ones, I think. 

\textbf{Lucas:} Yeah. 

\textbf{Lucas:} Um. So this question might not be fully applicable. When we were designing it, our primary users we had in mind were blind users, because they would obviously need the voice control as well. \textit{[P3: Yeah.]} Um, so when you're doing a task, how does your computer let you know what's happening? What it's doing? Um, the intent behind the question for blind users is obviously they can't see this. 

\textbf{P3:} Right, I totally understand. Yeah. Uh, it would be a real challenge for a blind user, because, like, I give the commands, and I see them on the screen. Whether it be my telephone, the TV, the computer. I can visually see. And I know that I know it'd move on. \textit{[Lucas: Yeah.]} To whatever, right? Or if you were blind, it doesn't tell you anything. \textit{[Lucas: Yeah.]} It doesn't say "task completed" or anything. 

\textbf{Lucas:} Okay. So Voice Access, what you use currently, you have to be able to see the screen. 

\textbf{P3:} Yeah, yeah, it would not work for someone who was visually impaired. 

\textbf{Lucas:} Gotcha. Um, yeah, no, that's good enough for there. Um, skip that one. So do you ever have it that when you're using Voice Access or any sort of accessibility tool, does it ever do anything unexpected? Like, suddenly goes off on its own, and how do you remedy that situation? 

\textbf{P3:} Yeah. That's, uh, when it misinterprets what you say, and you could say, uh, like the number six. \textit{[Lucas: Yeah.]} That's what I have trouble with all the time. I don't know if it's my pronunciation or what, I imagine most people are under the, I say "six," and that's things I'm saying sex. It comes up with some funny ones, though, but that's just one. \textit{[Lucas: Yeah.]} One thing that I ran into all the time. Um, there are other, that, uh, besides the Voice Access I use, Google has their own version, uh, that will help you. The telephone, I know this, so... It's Google. \textit{[Lucas: Yeah.]} I have Google and Voice Access for the computer and it's a different one for the TV. I can't remember what it uses, which, uh, program it uses, but there are other options out there. 

\textbf{Lucas:} Yeah. So if you say you had told it to copy the contents of one cell to another, and say copy it to the wrong cell. Yeah. Do you have to manually then say, uh, copy it back to the old cell, or can you just tell it, like, undo the last action or undo the last five minutes of that? 

\textbf{P3:} Yeah, you can. And if you're lucky, it will do that. 

\textbf{Lucas:} Okay. 

\textbf{P3:} If you're not lucky, you gotta go back to fixing it by hand. It's terrible, to go back to the very beginning and do everything over again. So very time consuming, it can be. 

\textbf{Lucas:} So, the functionality where you can just say "undo this," it doesn't work too well then. 

\textbf{P3:} No, no. It will for simple things, like if you're just copying, pasting. You know, it'll undo stuff like that. But if it's more complex, what you're doing, it's just lost. It's gone. 

\textbf{Lucas:} Gotcha. 

\textbf{P3:} So, yeah. 

\textbf{Lucas:} Um... So, you had mentioned at the start that, now that you've lost functionality of your arms to control the computer, there are some tasks that you can't quite do as well. So is there anything on the computer that you can't currently do or that you used to be able to do, that you would like to do? Or you say, hey, I think they should be able to add this functionality. 

\textbf{P3:} It's just, it's, as great as it is to be able to use your voice, it's very concerning that you may say something incorrect, or that it misinterprets what you mean, and it just goes ahead and does whatever task it thinks you mean. And if you don't have a backup, then you just, you're out of luck, it's gone. 

\textbf{Lucas:} Yes. 

\textbf{P3:} Yeah, you can lose important documentation or, uh, even a whole program, you could, you could, uh, cause damage that you wouldn't get back. 

\textbf{Lucas:} So when you're using Voice Access, does it ever ask you for permission? Like, does it ask you to confirm your actions before you do that? 

\textbf{P3:} Yeah, sometimes, yeah. Sometimes. Not always, so. 

\textbf{Lucas:} So do you want it? Because I know it's like phones and stuff, two-factor authentication. Like, I appreciate the security, but sometimes it seems too much. Would you rather it confirm lots or? 

\textbf{P3:} Yeah. Like, uh, for instance, when I'm doing my banking, that some banks will send you, just when you're getting on, that's verifying, that's actually you. \textit{[Lucas: Yeah.]} You've got passwords, I mean, you log in, you give it all that. Then it wants further proof that it's you, not a hacker. \textit{[Lucas: Yeah.]} And so, it goes by voice, and usually that's pretty good, but sometimes it misinterprets what you say or whatever. Uh, and that's definitely an issue. 

\textbf{Lucas:} Um, so say you were doing some more complex operation, a spreadsheet, where you realize that if this messes up, you're going to have to start from scratch. Would you want it to ask for confirmation then? Or do you feel like having every request say, "yes, go ahead," is cumbersome? 

\textbf{P3:} Yeah, yeah, its good. The one thing, like right now, uh, it definitely takes your voice as one of the proofs of identity that you are. And I was... Oh, the one I have to, uh, it'll send me a six-digit code to verify to my email or to my cell phone. Well, cell phone. It doesn't always work for me. \textit{[Lucas: Yeah.]} With the voice. So it's easier for me to have the six or nine digit number sent to my email address. Whereas CIBC won't do that. They don't figure that secure enough. \textit{[Lucas: Yeah.]} So that you have more tasks to do with them, which sometimes is easy and sometimes not. 

\textbf{Lucas:} Yeah. So is your phone hooked up to your computer? Are they totally separate? 

\textbf{P3:} Uh... No, it's strange because I can be talking on my cell phone. \textit{[Lucas: Yeah.]} Which is the same telephone number as my computer, yet I can have two conversations going on. 

\textbf{Lucas:} Really? 

\textbf{P3:} At the same number. Yeah, I know. I don't understand quite how that works, but... 

\textbf{Lucas:} Huh. Yeah. So, there, it's not really like a seamless connection between the two. 

\textbf{P3:} No. 

\textbf{Lucas:} Yeah, it's not like if you were on your computer, you can just say, copy the code from my phone onto there. 

\textbf{P3:} Yeah, yeah. They'll get there, but they're not there yet, at least not that I'm aware of. So, that answered your question about the security. 

\textbf{Lucas:} Yep. 

\textbf{P3:} Sometimes, like, it's overkill. \textit{[Lucas: Yeah.]} And then sometimes you're glad it is. \textit{[Lucas: Yes.]} Like, as much of a pain as it is, that extra step with CIBC, it's a pain at the time, but really, I appreciate the security. \textit{[Lucas: Yeah.]} It's tougher for, like, somebody to copy my voice, right? So, yeah, that's definitely something important with voice. 

\textbf{Lucas:} Um, so you've mentioned a few times in passing, kind of, that when they're asking for authentication, they're using your voice. Where is that? What's that being used? 

\textbf{P3:} Uh, banking, some of that you're giving in there to pay bills and everything. \textit{[Lucas: Okay. Yeah.]} There's a... Banking, uh, one's called Investor's Edge, where you can purchase GICs, buy stocks, any of the regular everyday banking features that, uh, yeah, have to do. 

\textbf{Lucas:} Yeah. So does that, to your knowledge, does the checking if it's you through the voice, is that from the website? Like, does the website allow you to set that up, or is that through your Voice Access that's checking it? 

\textbf{P3:} That's Google, that's Voice Access. That's the program. 

\textbf{Lucas:} Okay, so it's kind of the tooling. Make sure, make sure that you... 

\textbf{P3:} Yeah, it's mainly. It's, um, but it's, uh, mainly the financial institutions themselves, whether it be BMO or whatever. \textit{[Lucas: Yeah.]} They, they always want, they use your voice. 

\textbf{Lucas:} Interesting. Because, like, when I'm logging in to my banking, usually they'll, yeah, like, give it the code to an email or your phone number, something. I can't say that I've seen, like, set up a voice command. Yeah. Interesting, interesting. I don't know 

\textbf{P3:} I don't know how it tells the difference. 

\textbf{Lucas:} Yeah, yeah. But you do use your voice. 

\textbf{P3:} Yeah, it must give you the options. At some point. 

\textbf{Lucas:} Um, is there anything else that you would like to share that would be important to us about your understanding of how to use a computer? Is there anything that we haven't covered thus far? 

\textbf{P3:} Don't really think so. Not that I've, I'm sure there's way more, the more complex the task, the more difficult it would be. \textit{[Lucas: Yeah.]} And I'm sure there are some things that won't do. But for basic, uh, everyday word processing, basic spreadsheets, you can get along pretty well. 

\textbf{Lucas:} Yeah. Is there anything, so we've kind of talked a decent amount about what the pain points are. Is there anything that the existing stuff does really well? 

\textbf{P3:} Yeah, a lot of the everyday tasks, like the word processing, is great other than it misinterpreting what you're saying. \textit{[Lucas: Okay. Yeah.]} And even the basic spreadsheets, I use the most spreadsheets. \textit{[Lucas: Yeah.]} They're pretty, pretty easy to use and work, work actually. Other than misinterpreting what you're saying. 

\textbf{Lucas:} Yeah, yeah. So if you had a perfect input, it would do exactly what you're saying. 

\textbf{P3:} Yeah, correct. 

\textbf{Lucas:} Okay, let's see here... 

\textbf{P3:} And one other thing about the voice... \textit{[Lucas: Yeah.]} It's not as accurate. It's not as precise, like that LipSync I've got there. That's precise. I can make sure I get the exact wording, the exact everything, numbers. Whereas the Voice Access, it's still not as accurate. It's not as right. 

\textbf{Lucas:} So, how does the LipSync work? Like, what is LipSync? 

\textbf{P3:} Okay. I'll show you. If you just come over here for a sec. 

\textbf{Lucas:} Yes. 

\textbf{P3:} Okay, see that, uh... That mouthpiece. 

\textbf{Lucas:} Yep. 

\textbf{P3 (demonstration):} Can you take your hand. See the pointer on the screen. It's upper right. That piece there, I pull up, and I get my mouth around that. And you just use your fingers. You can take all of that and watch it move the mouse. Yeah, and it can move all around the screen. And then, uh, to do the various tasks, you can either breathe in, or blow out. \textit{[Lucas: Yeah.]} And you can move it left to right. \textit{[Lucas: Okay.]} Up and down, and that's my, that's my mouse. \textit{[Lucas: Yeah.]} And that's way more accurate than voice at this point. 

\textbf{Lucas:} Yeah. But I imagine it would be much slower. 

\textbf{P3:} It is slower, yeah, definitely. It is slower, but it's, it's more accurate. Like banking, well, I have to use that because it won't let me use voice. 

\textbf{Lucas:} Gotcha. Banking, yeah. 

\textbf{P3:} Yeah. 

\textbf{Lucas:} Okay, that's interesting. 

\textbf{P3:} And that's what I believe is Canadian and it's called LipSync. 

\textbf{Lucas:} LipSync. Okay, um, let me just make sure I got everything here. Yeah, I think that's all the major questions I've got. If you've got nothing else that you want... 

\textbf{P3:} Yeah, no, not really. Just the accuracy and, uh, so, at this point in time, they don't have, you can't do any banking with your voice. So banking's a big one. 

\textbf{Lucas:} It's a big one. Yeah, exactly. 

\textbf{P3:} Even for getting in, for someone, uh, like who was blind, whatever, like, you just couldn't do any banking at all. 

\textbf{Lucas:} Yeah, perfect. Well, thanks for having me. 

\textbf{P3:} Absolutely. If there's something, uh, more that you need in the future with this or anything else, just let me know. 

\textbf{Lucas:} That was going to be the last question, is we may need to do usability testing. Would you be interested in that? 

\textbf{P3:} Absolutely, yeah. 

\textbf{Lucas:} Perfect. Well, I will reach out to you then. 

\bigskip
\hrule
\bigskip

\textbf{Questions not covered:}
\begin{itemize}
  \item Q9 was skipped as it did not apply to this user based on his response to Q8.
\end{itemize}
